\documentclass{article}
\usepackage{amsmath}
\usepackage{enumerate}

% ==============================================================================
% VIDYAHEIST CONTENT TEAM GUIDE: BIOLOGY MCQ FORMAT
% ==============================================================================
% 
% 1. MAIN LIST: Wrap all questions in a \begin{enumerate} ... \end{enumerate}
% 2. EACH QUESTION: Start with \item followed by the question text.
% 3. ANSWERS: Add \hfill \textbf{(A)} (or B, C, D) at the very end of the question text line.
% 4. OPTIONS: 
%    - Option A: Standard nested \begin{enumerate}[(A)] ... \end{enumerate}
%    - Option B: Inline format like (A) Option A (B) Option B etc.
% ==============================================================================

\begin{document}

\title{Blueprint: Biology Question Paper}
\author{Vidyaheist Content Team}
\date{}
\maketitle

\begin{enumerate}

    % --- Standard MCQ Example ---
    % Tip: \hfill \textbf{(B)} tells the website B is the correct answer.
    \item Which cell organelle is known as the "Powerhouse of the cell"? \hfill \textbf{(B)}
    \begin{enumerate}[(A)]
        \item Nucleus
        \item Mitochondria
        \item Ribosome
        \item Lysosome
    \end{enumerate}

    % --- Match the Following Example ---
    % Tip: You can put descriptions on new lines using \\
    \item Match the following parts of a flower with their functions: \hfill \textbf{(A)}
    \\1. Stamen $\rightarrow$ (a) Male reproductive part
    \\2. Pistil $\rightarrow$ (b) Female reproductive part
    \\3. Petals $\rightarrow$ (c) Attract insects
    \\4. Sepals $\rightarrow$ (d) Protect the bud
    \\(A) 1-a, 2-b, 3-c, 4-d
    \\(B) 1-b, 2-a, 3-d, 4-c
    \\(C) 1-c, 2-d, 3-a, 4-b
    \\(D) 1-d, 2-c, 3-b, 4-a

\end{enumerate}

\end{document}
